\section{Methodological Formulation and Execution Pipeline}

\subsection{Algorithmic Overview}
The \texttt{mambashedder} module executes a structured transformation replacing monolithic dense projections with mathematically optimized submodules. The pipeline operates in five consecutive phases:
\begin{enumerate}
\item \textbf{Conditioning Analysis:} Verify matrix rank, Frobenius norm, and condition number.
\item \textbf{Decomposition / Sparsification:} Apply SVD spectral factoring or proximal shrinkage.
\item \textbf{Constraint Projection:} Enforce rank truncation via cumulative energy $\tau$ or 2:4 tensor-core masks.
\item \textbf{Residual Adapter Injection:} Attach zero-initialized LoRA branches.
\item \textbf{Safety Fallback Audit:} Revert to uncompressed weights if relative reconstruction error breaches safety limits.
\end{enumerate}

\begin{figure*}[t]
\centering
\begin{verbatim}
+-------------------------------------------------------------------------------+
|                      Pretrained Weight Matrix W                               |
+---------------------------------------+---------------------------------------+
                                        |
                                        v
+-------------------------------------------------------------------------------+
|  Phase 1: SVD / Proximal Factorization: W = U Sigma V^T OR S\_lambda(W)        |
|  Phase 2: Energy Rank / Structured Mask: E(r) \textgreater{}= tau OR Pi\_{2:4}(W)           |
|  Phase 3: Symmetric Splitting: U\_r = U Sigma^{1/2}, V\_r = Sigma^{1/2} V^T     |
|  Phase 4: Adapter Binding: W\_eff = U\_r V\_r + (alpha / r) B A (B init to 0)    |
|  Phase 5: Fallback Guardrail: If ||W - W\_eff||\_F / ||W||\_F \textgreater{} max\_err -\textgreater{} Dense |
+-------------------------------------------------------------------------------+
\end{verbatim}
\caption{End-to-end execution and safety verification pipeline of \texttt{mambashedder}.}
\end{figure*}

\subsection{Formal Mathematical Derivations}

\subsubsection{Linear Projection Mapping}
Let $X \in \mathbb{R}^{\mathcal{B} \times N \times d_{\text{in}}}$ denote activation inputs. Standard linear projection computes:
\[
Z = X W^\top + b
\]

\subsubsection{Low-Rank Bilinear Factorization}
Under low-rank compression, weight matrix $W \in \mathbb{R}^{d_{\text{out}} \times d_{\text{in}}}$ is factored into matrices $V_r \in \mathbb{R}^{r \times d_{\text{in}}}$ and $U_r \in \mathbb{R}^{d_{\text{out}} \times r}$:
\[
\hat{W} = U\_r V\_r
\]

\subsubsection{Factorized Forward Execution}
The forward evaluation transforms from monolithic matrix multiplication into two lightweight steps:
\[
\tilde{Z} = (X V\_r^\top) U\_r^\top + b
\]

\subsubsection{Cumulative Spectral Energy}
Rank selection is governed by cumulative squared singular values:
\[
\mathcal{E}(r) = \frac{\sum\_{i=1}^{r} \sigma\_i^2}{\sum\_{j=1}^{\min(d\_{\text{in}}, d\_{\text{out}})} \sigma\_j^2}
\]

\subsubsection{Energy-Thresholded Optimal Rank}
For target energy threshold $\tau \in (0, 1)$, the optimal rank $r^*$ satisfies:
\[
r^* = \min \left\{ r \in \mathbb{N} : \mathcal{E}(r) \ge \tau \right\}
\]

\subsubsection{Soft-Thresholding Shrinkage}
In proximal sparsity modes, the $\ell_1$ proximal operator executes element-wise:
\[
\mathcal{S}\_\lambda(W) = \operatorname{sign}(W) \odot \max(|W| - \lambda, 0)
\]

\subsubsection{Hardware-Aligned 2:4 Structured Sparsity}
For NVIDIA tensor core execution, each contiguous 4-element block is projected to 50\% sparsity:
\[
\Pi\_{2:4}(W\_{i, 4k:4k+4}) = W\_{i, 4k:4k+4} \odot M\_{i, 4k:4k+4}, \quad \sum\_{j=0}^{3} M\_{i, 4k+j} = 2
\]

\subsubsection{Iterative Proximal Update}
Calibration updates parameters via projected gradient steps:
\[
W^{(t+1)} = \Pi\_{\Omega}\left( W^{(t)} - \eta \nabla f(W^{(t)}) \right)
\]

\subsubsection{Residual Adapter Injection}
To recover expressive capacity lost during aggressive compression, trainable low-rank adapters are attached:
\[
Y = X \hat{W}^\top + \frac{\alpha}{r\_{\text{lora}}} (X A^\top) B^\top + b
\]

\subsubsection{Zero-Residual Initialization}
Initial equivalence with unadapted weights is guaranteed by zero initialization:
\[
A \sim \mathcal{N}\left(0, \frac{1}{r\_{\text{lora}}}\right), \quad B = 0 \implies B A = 0
\]

\subsubsection{Relative Frobenius Reconstruction Error}
Post-compression parameter distortion is evaluated against the uncompressed matrix:
\[
\delta\_{\text{rel}}(W, \hat{W}) = \frac{\|W - \hat{W}\|\_F}{\|W\|\_F + \epsilon}
\]

\subsubsection{Automated Guardrail Decision Rule}
Should the relative error breach pre-specified threshold $\delta_{\text{max}}$:
\[
\operatorname{Module}(W) = \begin{cases} \hat{W} + \frac{\alpha}{r} B A & \text{if } \delta\_{\text{rel}} \le \delta\_{\text{max}} \text{ and } \operatorname{isfinite}(\hat{W}) \\ W & \text{otherwise (safety fallback)} \end{cases}
\]

\subsubsection{Attention Kernel Normalization}
In linear attention formulations, outputs are computed via associative reduction:
\[
O\_t = \frac{\sum\_{s=1}^{t} \exp(-\gamma (t-s)) \phi(Q\_t) (\phi(K\_s)^\top V\_s)}{\sum\_{s=1}^{t} \exp(-\gamma (t-s)) \phi(Q\_t) \phi(K\_s)^\top + \epsilon}
\]

\subsubsection{Head Projection Partitioning}
Multi-head attention projections decompose into parallel heads:
\[
W\_Q = [W\_Q^{(1)}; W\_Q^{(2)}; \dots; W\_Q^{(H)}], \quad W\_Q^{(h)} \in \mathbb{R}^{d \times d\_{\text{in}}}
\]

\subsection{Detailed Algorithmic Execution Trace}

\textbf{Algorithm 1:} Complete Structured Execution of \texttt{MambaShedderCompressor}
\begin{lstlisting}[language=Python]
# Phase 1: Conditioning and norm evaluation
W\_f = layer.weight.data.float()
frob\_norm = torch.linalg.norm(W\_f)
assert frob\_norm \textgreater{}= 1e-4, 'Degenerate weight norm'

# Phase 2: Core mathematical decomposition
if is\_svd\_factorization:
    U, S, Vh = torch.linalg.svd(W\_f, full\_matrices=False)
    cum\_energy = torch.cumsum(S ** 2, dim=0) / torch.sum(S ** 2)
    r = torch.searchsorted(cum\_energy, tau).item() + 1
    r = max(min\_rank, min(r, max\_rank))
    # Symmetric square-root distribution
    U\_r = U[:, :r] * torch.sqrt(S[:r])
    V\_r = torch.sqrt(S[:r]).unsqueeze(1) * Vh[:r, :]
elif is\_proximal\_sparsity:
    W\_curr = W\_f.clone()
    for step in range(num\_steps):
        # Soft-thresholding
        W\_soft = torch.sign(W\_curr) * torch.clamp(torch.abs(W\_curr) - lambda\_prox, min=0.0)
        # Structural 2:4 or block projection
        W\_proj = apply\_structural\_mask(W\_soft, sparsity\_target)
        W\_curr = W\_curr - eta * (W\_curr - W\_proj)
    W\_r = W\_proj

# Phase 3: Adapter binding & zero-initialization
lora\_A = nn.Parameter(torch.randn(lora\_r, d\_in) / math.sqrt(lora\_r))
lora\_B = nn.Parameter(torch.zeros(d\_out, lora\_r))

# Phase 4: Fallback verification
recon\_err = torch.linalg.norm(W\_f - W\_reconstructed) / frob\_norm
if recon\_err \textgreater{} max\_error\_ceiling or torch.isnan(recon\_err):
    logger.warning('Triggering safety fallback to uncompressed layer')
    layer = uncompressed\_layer\_backup
\end{lstlisting}

\subsection{Complexity and FLOPs Analysis}
\subsubsection{Computational Complexity Breakdown}
Evaluating a dense linear projection of dimension $d_{\text{in}} \times d_{\text{out}}$ requires $2 d_{\text{in}} d_{\text{out}}$ FLOPs per sequence token. For low-rank compressed factorization at rank $r$, the sequential forward evaluation requires:
\[
\text{FLOPs}\_{\text{factorized}} = 2 N d\_{\text{in}} r + 2 N r d\_{\text{out}} = 2 N r (d\_{\text{in}} + d\_{\text{out}})
\]
The theoretical FLOP reduction ratio $\rho_{\text{FLOPs}}$ is therefore:
\[
\rho\_{\text{FLOPs}} = \frac{2 N d\_{\text{in}} d\_{\text{out}}}{2 N r (d\_{\text{in}} + d\_{\text{out}})} = \frac{d\_{\text{in}} d\_{\text{out}}}{r (d\_{\text{in}} + d\_{\text{out}})}
\]
For symmetric hidden projections with $d_{\text{in}} = d_{\text{out}} = d$:
\[
\rho\_{\text{FLOPs}} = \frac{d^2}{2 r d} = \frac{d}{2 r}
\]
When $d = 768$ and $r = 128$, the mathematical speedup reaches $3.0	imes$.

\begin{table}[H]
\centering
\small
\begin{tabular}{lcccc}
\toprule
\textbf{Operation} & \textbf{Standard} & \textbf{mamba\_shedder} & \textbf{Complexity} & \textbf{FLOPs ($d=768$)} \\
\midrule
Q Projection & $d \times d$ & $2 \times (d \times r)$ & $\mathcal{O}(2rd)$ & $0.196$ M \\
K Projection & $d \times d$ & $2 \times (d \times r)$ & $\mathcal{O}(2rd)$ & $0.196$ M \\
V Projection & $d \times d$ & $2 \times (d \times r)$ & $\mathcal{O}(2rd)$ & $0.196$ M \\
Output Proj & $d \times d$ & $2 \times (d \times r)$ & $\mathcal{O}(2rd)$ & $0.196$ M \\
LoRA Residual & None & $2 \times (d \times r_L)$ & $\mathcal{O}(2r_L d)$ & $0.024$ M \\
\midrule
\textbf{Total Attention Proj} & $4 d^2$ & $8rd + 8r_L d$ & $\mathcal{O}(rd)$ & \textbf{0.808 M (vs 2.36 M)} \\
\bottomrule
\end{tabular}
\caption{Computational FLOPs breakdown for multi-head attention projections.}
\end{table}

\subsection{Gradient Propagation and Backpropagation Dynamics}
Computing gradients through decomposed low-rank projections requires careful backward chaining. Let $\delta Y = \frac{\partial \mathcal{L}}{\partial Y} \in \mathbb{R}^{\mathcal{B} \times N \times d_{\text{out}}}$ denote the incoming upstream gradient. Under factorized linear projection $\hat{Y} = X V_r^\top U_r^\top$, the gradients with respect to factor matrices $U_r$ and $V_r$ are derived as:
\begin{align}
\frac{\partial \mathcal{L}}{\partial U\_r} &= \sum\_{b=1}^{\mathcal{B}} (\delta Y\_b)^\top (X\_b V\_r^\top) \in \mathbb{R}^{d\_{\text{out}} \times r} \\
\frac{\partial \mathcal{L}}{\partial V\_r} &= \sum\_{b=1}^{\mathcal{B}} ((\delta Y\_b) U\_r)^\top X\_b \in \mathbb{R}^{r \times d\_{\text{in}}} \\
\frac{\partial \mathcal{L}}{\partial p{5cm}} &= (\delta Y) U\_r V\_r \in \mathbb{R}^{\mathcal{B} \times N \times d\_{\text{in}}}
\end{align}
When factor matrices are kept frozen during inference recovery, computing gradient caches for base factors is bypassed entirely, cutting backward memory by up to $65\%$.

\subsection{Hardware-Software Co-Design and Memory Footprint Mapping}
Modern deep learning accelerators (e.g. NVIDIA A100, H100) feature specialized matrix multiplication units (Tensor Cores) and hierarchical on-chip memories. In a standard uncompressed linear projection $Y = X W^\top$, the processor loads the full weight tensor $W \in \mathbb{R}^{d_{\text{out}} \times d_{\text{in}}}$ from High Bandwidth Memory (HBM) into streaming multiprocessor (SM) register files. Under factorized execution, the matrix multiplication is partitioned into two sequential operations:
\begin{align}
H &= X V\_r^\top \quad \text{with FLOPs } 2 \mathcal{B} N d\_{\text{in}} r \\
Y &= H U\_r^\top \quad \text{with FLOPs } 2 \mathcal{B} N r d\_{\text{out}}
\end{align}
When rank $r$ satisfies $r < \frac{d_{\text{in}} d_{\text{out}}}{d_{\text{in}} + d_{\text{out}}}$, total arithmetic operations are strictly reduced. Furthermore, when intermediate hidden activations $H$ are kept in fast on-chip SRAM (L1/shared memory), off-chip memory traffic is reduced by the ratio $\frac{d_{\text{in}} d_{\text{out}}}{r(d_{\text{in}} + d_{\text{out}})}$, directly circumventing the memory wall.

\begin{figure}[H]
\centering
\begin{verbatim}
+-------------------------------------------------------------+
|               Tensor Memory Pipeline (On-Chip SRAM)         |
+-------------------------------------------------------------+
|  Global HBM DRAM                                            |
|    |                                                        |
|    |-- Load X [B, N, D\_in] ----------------+                |
|    |-- Load V\_r [r, D\_in] --------+        |                |
|                                   |        |                |
|                                   v        v                |
|  Shared Memory / Registers   [ Compute H = X V\_r^T ]        |
|                                   |                         |
|    |-- Load U\_r [D\_out, r] -------|                         |
|                                   v                         |
|  Shared Memory / Registers   [ Compute Y = H U\_r^T ]        |
|                                   |                         |
|                                   v                         |
|  Global HBM DRAM             [ Store Output Y ]             |
+-------------------------------------------------------------+
\end{verbatim}
\caption{Hardware-aligned memory hierarchy and on-chip SRAM data reuse.}
\label{fig:memory\_hierarchy}
\end{figure}

\subsection{Rigorous Proof of Theoretical Reconstruction Bounds}

\textbf{Theorem 1 (Optimal Energy Truncation Bound).} \textit{Let $W \in \mathbb{R}^{m \times n}$ have singular values $\sigma_1 \ge \sigma_2 \ge \dots \ge \sigma_{\rho} > 0$. If rank $r$ is selected such that $\mathcal{E}(r) \ge \tau$, then the relative Frobenius reconstruction error is bounded strictly by $\sqrt{1 - \tau}$.}

\begin{proof}
By the singular value decomposition, the squared Frobenius norm of $W$ equals the sum of its squared singular values:
\[
\|W\|\_F^2 = \operatorname{Tr}(W^\top W) = \sum\_{i=1}^{\rho} \sigma\_i^2
\]
By construction, the truncated matrix $W_r = U_r V_r$ truncates all components with index $i > r$. Thus, the residual error matrix $\Delta W = W - W_r$ has singular values $\sigma_{r+1}, \dots, \sigma_{\rho}$.
Its squared Frobenius norm is:
\[
\|W - W\_r\|\_F^2 = \sum\_{i=r+1}^{\rho} \sigma\_i^2 = \sum\_{i=1}^{\rho} \sigma\_i^2 - \sum\_{i=1}^{r} \sigma\_i^2
\]
Dividing both sides by $\|W\|_F^2$:
\[
\frac{\|W - W\_r\|\_F^2}{\|W\|\_F^2} = 1 - \frac{\sum\_{i=1}^{r} \sigma\_i^2}{\sum\_{i=1}^{\rho} \sigma\_i^2} = 1 - \mathcal{E}(r)
\]
Since $r$ is chosen such that $\mathcal{E}(r) \ge \tau$, we have:
\[
\frac{\|W - W\_r\|\_F^2}{\|W\|\_F^2} \le 1 - \tau \implies \frac{\|W - W\_r\|\_F}{\|W\|\_F} \le \sqrt{1 - \tau}
\]
This completes the proof. $\square$
\end{proof}

\textbf{Theorem 2 (Initial Inference Invariance under Zero Residual).} \textit{Let the effective weight matrix be parameterised as $W_{\text{eff}} = U_r V_r + \frac{\alpha}{r_L} B A$. If adapter up-projection $B$ is initialized to zero ($B = 0$), then $W_{\text{eff}}$ evaluated at initialization step $t=0$ identically equals the truncated base projection $W_r$.}

\begin{proof}
At initialization $t=0$, $B = 0 \in \mathbb{R}^{d_{\text{out}} \times r_L}$. For any arbitrary matrix $A \in \mathbb{R}^{r_L \times d_{\text{in}}}$:
\[
B A = 0 \cdot A = 0 \in \mathbb{R}^{d\_{\text{out}} \times d\_{\text{in}}}
\]
Consequently:
\[
W\_{\text{eff}} = U\_r V\_r + \frac{\alpha}{r\_L} (0) = U\_r V\_r = W\_r
\]
Hence, output activations satisfy:
\[
Y = X W\_{\text{eff}}^\top + b = X W\_r^\top + b
\]
The initial forward pass preserves the exact mathematical characteristics of the truncated base projection without noise introduced by uncalibrated random adapter weights. $\square$
\end{proof}
